\chapter*{Resumo}
A precisão alcançada pelas medições realizadas no Large Hadron Collider (LHC) impõe
exigências cada vez mais rigorosas às previsões teóricas. Em Cromodinâmica Quântica
(QCD) perturbativa, a precisão destas previsões é sistematicamente melhorada através
da inclusão de correções de ordem superior na constante de acoplamento forte,
\(\alpha_s\). Contudo, aumentar a precisão de um cálculo não corresponde apenas a
acrescentar diagramas de Feynman: exige controlar uma estrutura complexa de
contribuições de emissão real e virtual, bem como as divergências que surgem em cada
uma delas.

As divergências infravermelhas constituem uma das principais dificuldades dos 
cálculos perturbativos em QCD. As contribuições de emissão real tornam-se singulares 
quando um partão é emitido com energia muito baixa, o limite \textit{soft}, ou quando dois
partões se tornam colineares. Estas singularidades estão associadas à emissão de
radiação não resolvida, isto é, de partões cuja energia ou separação angular é
demasiado pequena para serem distinguidos como estados finais distintos. Paralelamente,
as amplitudes que incluem correções virtuais de loop apresentam polos
infravermelhos associados às mesmas regiões cinemáticas. Embora estas divergências se 
cancelem em observáveis físicos suficientemente inclusivos, as componentes 
individuais do cálculo não são finitas. É, por isso, necessário reorganizar a 
previsão teórica de modo a isolar as singularidades e a efetuar o seu cancelamento 
de forma analítica e controlada.

O formalismo de subtração por antenas fornece um enquadramento adequado para esse
objetivo. Neste contexto, uma antena é formada por dois partões ligados por cor,
designados por radiadores, entre os quais ocorre a emissão de radiação. A radiação
não resolvida é então descrita em termos destas antenas, permitindo tratar de forma
sistemática as regiões singulares do espaço de fases. As funções de antena reproduzem os
limites \textit{soft} e colineares relevantes de forma local no espaço de fase. Quando são 
subtraídas às contribuições reais, eliminam as singularidades antes da integração 
numérica. Depois de integradas analiticamente sobre o espaço de fase da radiação 
não resolvida, originam polos em regularização dimensional que cancelam os polos 
das contribuições virtuais. O cancelamento destas singularidades permite que cada 
contribuição seja integrada numericamente em quatro dimensões.

Esta dissertação apresenta o desenvolvimento de \textit{AntCalc}, um pacote simbólico para a 
construção e integração automatizadas de funções de antena em QCD. O projeto 
procura automatizar etapas que requerem implementações específicas para cada antena, 
reunindo num mesmo fluxo as convenções de normalização, construção e integração. O
trabalho centra-se nas antenas
quark–antiquark sem massa, considerando a configuração cinemática final-final, na
qual ambos os radiadores da antena são partículas do estado final, até à ordem
next-to-next-to-leading order (NNLO). Para além deste objetivo, é explorada
uma extensão ao regime de quarks massivos, através da implementação da antena
\(A_3^0(Q,g,\bar{Q})\), que descreve a emissão de um gluão entre um quark e um
antiquark massivos.

A base teórica da dissertação inclui uma introdução à QCD perturbativa, aos campos de 
quarks e gluões e à construção de amplitudes a partir de diagramas de Feynman. São
discutidas as divergências ultravioleta, tratadas através da renormalização, e as
divergências infravermelhas, que se cancelam mediante a combinação das diferentes
contribuições perturbativas reais e virtuais à mesma ordem. Utiliza-se regularização
dimensional, na qual o número de dimensões
do espaço-tempo é tomado como \(d=4-2\epsilon\). Neste enquadramento, as divergências
manifestam-se como polos em \(1/\epsilon\), e as quantidades integradas são expressas por
séries de Laurent. O cancelamento dos polos pode, assim, ser acompanhado explicitamente 
ao combinar as diferentes contribuições para uma quantidade física.

É também apresentada a estrutura das funções de antena e a sua decomposição em cor. 
As antenas são obtidas a partir de amplitudes físicas normalizadas e organizadas de 
acordo com os radiadores e as partículas emitidas. A fatorização do espaço de fase é 
essencial neste processo, pois permite separar o espaço de fase reduzido, associado 
aos radiadores após a emissão não resolvida ser removida, do espaço de fase 
correspondente à própria emissão. A integração pode então ser efetuada apenas sobre 
o espaço de fase da radiação não resolvida, independentemente do processo duro em
que a antena venha a ser inserida.

A integração analítica das funções de antena sobre o correspondente espaço de fases
de antena constitui uma das etapas mais exigentes na obtenção das antenas
integradas. Para realizar esta integração, as integrais de espaço de fases são
reescritas como integrais de loop através da técnica de reverse unitarity. Estas
são então reduzidas, através de identidades de integração por partes (IBP) e do
algoritmo de Laporta, a um conjunto finito de integrais mestres. Em casos adequados, recorre-se também à redução de
Passarino–Veltman. Uma vez identificados os integrais mestres e substituídas as 
respetivas expressões conhecidas, obtém-se a forma integrada da antena como uma 
expansão em \(\epsilon\). O resultado da redução fica, deste modo, expresso num 
conjunto limitado de integrais mestres, cujas expansões em \(\epsilon\) determinam 
as partes singular e finita da antena integrada.

O pacote \textit{AntCalc} foi implementado no ambiente \textit{Wolfram Mathematica} e integra 
ferramentas especializadas para as diferentes fases do cálculo. FeynArts é utilizado 
para gerar os diagramas de Feynman relevantes, enquanto \textit{FeynCalc} e \textit{FeynHelpers} apoiam 
a construção e a manipulação das amplitudes. A redução IBP é realizada com \textit{LiteRed2}. 
A arquitetura do pacote é orientada por perfis, isto é, por descrições estruturadas 
que especificam o processo, as amplitudes, as famílias de integrais, os fatores de 
normalização e os integrais mestres necessários. A informação específica de cada 
antena fica, assim, concentrada no respetivo perfil, enquanto as rotinas de 
construção e integração permanecem comuns aos diferentes casos.

O fluxo de trabalho divide-se em duas fases principais: construção e integração. 
Na fase de construção, o utilizador fornece a descrição do processo e o pacote gera 
os diagramas e amplitudes correspondentes. São depois efetuadas as manipulações 
algébricas necessárias para extrair a função de antena não integrada, ou as suas 
componentes de cor quando aplicável. Na fase de integração, essa expressão é 
convertida para uma forma apropriada à redução, associada à família de integrais 
correta e reduzida aos integrais mestres. Finalmente, a substituição dos resultados 
dos integrais mestres e a expansão em torno de \(\epsilon=0\) produzem a antena 
integrada. Ao longo deste fluxo são efetuados testes de consistência e apresentados 
diagnósticos do cálculo.

Como exemplo detalhado, é apresentada a derivação da antena \(A_3^0\), que descreve a
emissão de um gluão entre um quark e um antiquark sem massa a \textit{tree-level}
(onde \textit{tree-level} corresponde a contribuições de diagramas de Feynman sem
loops). Este
caso ilustra todo o percurso, desde a geração da amplitude até à expressão integrada.
A validação local da antena é efetuada através dos seus limites infravermelhos. No
limite \textit{soft}, verifica-se que a expressão reproduz o fator eikonal esperado; nos
limites colineares, confirma-se a recuperação da função de desdobramento
(\textit{splitting-function}) de Altarelli–Parisi apropriada. A concordância nestes limites constitui uma 
verificação local da expressão obtida, independente da comparação posterior da 
antena integrada com resultados da literatura.

Para além de \(A_3^0\), \textit{AntCalc} é aplicado a um conjunto de antenas sem massa até NNLO. 
São obtidas contribuições de dois, três e quatro partões, incluindo os conjuntos 
\(A_2^2\), \(A_3^1\) e \(A_4^0\), bem como antenas com outras estruturas de cor, tais como 
\(B_4^0\) e \(C_4^0\). Os resultados não integrados e integrados são comparados com 
resultados estabelecidos na literatura, verificando-se concordância. A concordância 
com os resultados publicados verifica, em simultâneo, as convenções de regularização 
dimensional e normalização, bem como o tratamento da estrutura de cor e do espaço de 
fase adotado na implementação.

A validação global é realizada através do cálculo do rácio hadrónico \(R\), no
limite de quarks sem massa, até NNLO em QCD, isto é, até \(\mathcal{O}(\alpha_s^2)\).
Este observável é definido como o rácio entre a secção eficaz inclusiva de produção
de hadrões em aniquilação \(e^+e^-\) e a secção eficaz para a produção de um par de
muões,
\[
R(s) = \frac{\sigma(e^+e^-\to\text{hadrões})}{\sigma(e^+e^-\to\mu^+\mu^-)}.
\]
Por
ser inclusivo, constitui um teste rigoroso da consistência do formalismo: os polos
infravermelhos provenientes das contribuições reais e virtuais devem cancelar
exatamente. O cálculo do \(R\)-ratio a partir dos objetos integrados produzidos por
\textit{AntCalc} confirma esse cancelamento e reproduz o resultado finito conhecido até NNLO.
A reprodução do \(R\)-ratio complementa, deste modo, os testes locais das antenas, 
verificando o cancelamento dos polos quando as diferentes contribuições são 
combinadas num observável físico.

A dissertação explora ainda a extensão ao caso massivo pela construção e integração 
da antena massiva \(A_3^0\). A presença de quarks massivos nesta antena altera a cinemática, a estrutura
do espaço de fase e os limites infravermelhos, eliminando, em particular, algumas 
singularidades colineares existentes no caso sem massa. Este cálculo fornece uma 
primeira implementação do fluxo massivo em \textit{AntCalc} e identifica as alterações 
necessárias à construção e integração de futuras antenas massivas.

Em síntese, o trabalho desenvolvido mostra que as etapas de construção, redução e
integração das funções de antena consideradas podem ser reunidas num único
procedimento automatizado. O pacote \textit{AntCalc} desenvolvido nesta tese reproduz
resultados conhecidos para antenas sem massa até NNLO, valida os limites
infravermelhos locais e recupera corretamente o \(R\)-ratio hadrónico. A
implementação obtida pode servir de base à inclusão de novas classes de antenas sem
exigir que todo o procedimento de construção e integração seja novamente implementado
de raíz.

\vspace{1cm}
\noindent\textbf{Palavras chave:} QCD; Subtração por Antenas; NNLO; Singularidades Infravermelhas; \textit{Mathematica}

\chapter*{Abstract}
The precision achieved by measurements at the Large Hadron Collider (LHC) places
increasingly stringent demands on theoretical predictions. In perturbative Quantum
Chromodynamics (QCD), the accuracy of these predictions is systematically improved
by including higher-order corrections in the strong coupling constant, \(\alpha_s\).
Beyond leading order, however, real and virtual contributions develop infrared
singularities associated with soft and collinear parton radiation. While these
singularities cancel in infrared-safe observables when all relevant contributions
are combined, the real and virtual terms are separately divergent and are defined
over phase spaces of different particle multiplicities. Their numerical evaluation
therefore requires a procedure that isolates the singular behaviour and implements
its cancellation explicitly.

Antenna subtraction provides a systematic way to do this. It uses universal antenna 
functions to describe unresolved radiation between colour-connected hard partons. 
These functions are subtracted locally from the real-emission contribution and then 
integrated analytically in dimensional regularisation, making the cancellation of
infrared poles against the virtual contributions explicit.

This thesis develops \textit{AntCalc}, a symbolic framework in Wolfram Mathematica 
for constructing and integrating QCD antenna functions. The package combines 
\textit{FeynArts}, \textit{FeynCalc}, \textit{FeynHelpers}, and \textit{LiteRed2}, 
while keeping process-dependent 
information separate from the general calculation workflow. It generates the required 
amplitudes, constructs unintegrated antennae or colour components, maps the 
corresponding phase-space integrals onto integral families, and reduces them to 
master integrals through integration-by-parts identities.

The implementation reproduces the massless final-final quark-antiquark antennae 
through next-to-next-to-leading order. The tree-level three-parton antenna \(A_3^0\) 
is used as a detailed example and is checked in its soft and collinear limits. The
integrated antenna functions are found to agree with existing results in the
literature, while their combined consistency is further validated through the
calculation of the massless hadronic \(R\)-ratio up to NNLO, \(\mathcal{O}(\alpha_s^2)\).
The extension beyond the massless case is also explored through the implementation
of the massive \(A_3^0\) antenna, demonstrating the applicability of the developed
framework to antennae involving massive partons. Future work will extend AntCalc to further massive antennae,
additional partonic configurations, and a more modular multi-language architecture.

\vspace{1cm}
\noindent\textbf{Keywords:} QCD; Antenna Subtraction; NNLO; Infrared Singularities; \textit{Mathematica}

\addcontentsline{toc}{chapter}{Resumo e Abstract}